\documentclass[10pt,a4paper]{amsart}
\usepackage{amsmath,amssymb,amsthm,graphicx}
%\usepackage{amsthm}
\usepackage{amsmath}
\usepackage{amssymb}
\usepackage{amsthm}
\usepackage{array}
\usepackage{xy}
\usepackage{hyperref}
\usepackage{color}
\usepackage{transparent}
\usepackage{latexsym} 
\theoremstyle{plain}
  \newtheorem{theorem}{Theorem}[section]
  \newtheorem{lemma}{Lemma}[section]
  \newtheorem{proposition}{Proposition}[section]
  \newtheorem{question}{Question}[section]
  \newtheorem{corollary}{Corollary}[section]
  \newtheorem{definition}{Definition}[section]
  \newtheorem{remark}{Remark}[section]
  \newtheorem{example}{Example}[section]
  \newtheorem{note}{Note}[section]
  \newtheorem{exercise}{Exercise}[section]
\newcommand{\y}{\bar{y}}
\newcommand{\U}{\bar{u}}
\newcommand{\z}{\bar{c}}

   \newcommand{\beqn}{\begin{eqnarray}}
   \newcommand{\eeqn}{\end{eqnarray}}
   \newcommand{\beqs}{\begin{eqnarray*}}
   \newcommand{\eeqs}{\end{eqnarray*}}
   \newcommand{\ban}{\begin{eqnarray*}}
   \newcommand{\nan}{\end{eqnarray*}}
   \newcommand{\beq}{\begin{equation}}
   \newcommand{\eeq}{\end{equation}}
   
   \newcommand{\p}{\partial}
\newcommand{\eps}{\varepsilon}
\newcommand{\Lm}{\Omega^*}
\newcommand{\Om}{\Omega}
\newcommand{\pom}{{\p\Om}}
\newcommand{\bom}{{\overline\Om}}

  \newcommand{\osc}{{\mbox{osc}}}
\renewcommand{\det}{\mbox{det}}
  \newcommand{\mes}{\mbox{mes}}
  \newcommand{\dist}{\mbox{dist}}
  \newcommand{\diam}{{\mbox{diam}}}
    \newcommand{\rad}{{\mbox{radius}}}
  \newcommand{\diag}{{\mbox{diag}}}
  \newcommand{\Vol}{{\mbox{Vol}}}
  \newcommand{\arctg}{{\mbox{arctg}}}
  \newcommand{\supp}{{\mbox{supp}}}

\newcommand{\ol}{\overline}
\newcommand{\ul}{\underline}
\newcommand{\lan}{\langle}
\newcommand{\ran}{\rangle}
\newcommand{\D}{\nabla}  

\newcommand{\s}{\tilde{v}}
\newcommand{\e}{\varepsilon}
\newcommand{\R}{\mathbb{R}}
\newcommand{\0}{\textbf{0}}
\numberwithin{equation}{section}
\DeclareMathOperator{\cexp}{c-exp}
\DeclareMathOperator{\cstarexp}{c*-exp}

%    Absolute value notation
\newcommand{\abs}[1]{\lvert#1\rvert}

%    Blank box placeholder for figures (to avoid requiring any
%    particular graphics capabilities for printing this document).
\newcommand{\blankbox}[2]{%
  \parbox{\columnwidth}{\centering
%    Set fboxsep to 0 so that the actual size of the box will match the
%    given measurements more closely.
    \setlength{\fboxsep}{0pt}%
    \fbox{\raisebox{0pt}[#2]{\hspace{#1}}}%
  }%
}


%%%%%%%%%%%%%%%%%%%%%%%
\renewcommand{\theequation}{\arabic{section}.\arabic{equation}}
%\newcommand{\mysection}[1]{\section{#1}\setcounter{equation}{0}}
  \numberwithin{equation}{section}
  \numberwithin{figure}{section}

%\TagsOnRight
%\newcommand{TagsOnRight}
%%%%%%%%%%%%%%%%%%%%%%%%%

\renewcommand{\baselinestretch}{1.00}
\parskip = 0.25in
\title{Holder differentiability of the separating singular hypersurface for transport into two convex targets without source convexity}
\author{Shibing Chen, Jiakun Liu}
\date{arxiv:2210.13841v4; CC BY 4.0}
\begin{document}
\maketitle
\noindent\textit{Editorial note.} Source599 overstates affine invariance for arbitrary distance-to-boundary divided by diameter. The actual bound602 is determined by the collinear section endpoints explicitly specified519 and the ratio574-579, which is affine-invariant. The original formula is retained; the independent review accepts this uniquely determined local explanation without altering the PDE argument.\par\medskip
\noindent\textit{Editorial note.} Source463 suppresses v(y\_k) in its affine-function display;451 states the correct value and v(y\_k) tends to v(p)=0, so the limiting argument is unchanged.\par\medskip
\noindent\textit{Editorial note.} The old pdftex graphicx driver is omitted from the wrapper preamble; the existing XeLaTeX graphicx driver imports the unchanged original diagram. No author mathematical text is edited.\par\medskip
\noindent\textit{Editorial scope.} One counted theorem t111, with all local supporting proofs353-678 as required solution. Standard prerequisites are only the named classical tools above; no local result is assumed. Requires independent substantive review before ready.\par\medskip
\noindent\textit{Allowed standard prerequisites (editorial).}\begin{itemize}
\item Brenier duality, interior Caffarelli regularity for convex targets, and the known Lipschitz separation of the two target regions. Use the transport potentials and initial Lipschitz hypersurface explicitly recalled in154-180 and303-316. No boundary Holder differentiability is given.. Br1,C92,CL1 Theorem5.1 and KM1 in retained author bibliography
\item John normalization and standard centred-section properties. Centred convex sections exist, are dimensionally balanced, depend continuously on positive height, and satisfy the barrier estimate339.. C92b; CM AppendixA, particularlyA.8
\item Localized Alexandrov estimates for Brenier solutions. Use la1/ua1 exactly as displayed365-374, with the upper estimate at a point whose subgradient meets an interior ball of the target. Constants may depend on the distance of that subgradient point to the target boundary. Also allow the same classical interior strict-convexity/differentiability result at a point inside the source whose subgradient meets the target interior, as invoked394 and632; this is Figalli-Kim Proposition3.4 in the author PDF (Proposition2 in the journal numbering).. Figalli-Kim2010, author version Proposition3.3 and equation3.1, PDF pages3-6; the paper cites final numbering Proposition1/equation4. Exact hypotheses read at https://people.math.ethz.ch/\textasciitilde{}afigalli/papers-pdf/Partial-regularity-of-Brenier-solutions-of-the-Monge-Ampere-equation.pdf
\item Classical localized boundary uniform convexity and convex duality. Once a convex support has two-sided Alexandrov density and uniformly localized centred sections near a boundary point, the local Caffarelli boundary method yields a power lower modulus for its gradient; Legendre duality gives Holder continuity of the inverse gradient.. C96 Lemma2.2/Corollary2.3 and CM Theorem7.13, as invoked665-669. Applicability must be assessed; no local differentiability or localization lemma of the counted article is supplied freely.
\end{itemize}\medskip
Let $f, g \in L^1(\mathbb{R}^n)$ represent two probability densities concentrated on bounded open sets $\Omega, \Omega^*\subset\mathbb{R}^n$, respectively. Suppose there exists a positive constant \(\lambda\) such that $\frac{1}{\lambda} < f, g < \lambda$ in $\Omega, \Omega^*$, respectively.
Invoking Brenier's theorem \cite{Br1}, we identify two globally Lipschitz convex functions, $u$ and $v$, both defined on $\mathbb{R}^n$, satisfying the following conditions:
\begin{align}
    (Du)_\sharp f = g, & \text{ with } Du(x) \in \overline{\Omega^*} \text{ for almost every } x \in \mathbb{R}^n; \label{condi1}\\
    (Dv)_\sharp g = f, & \text{ with } Dv(y) \in \overline{\Omega} \text{ for almost every } y \in \mathbb{R}^n. \label{condi2}
\end{align}
From the regularity theory of Caffarelli \cite{C92}, if $\Omega^*$ is convex, then the potential function $u$ belongs to $C^{1,\beta}_{\text{loc}}(\Omega)$ for some $\beta\in(0,1)$ and is strictly convex inside $\Omega$.

If $\Omega^*$ consists of two disjoint convex domains, as seen in the example of \cite{C92} when mapping the unit disc onto two shifted half-discs, there appears a singular set $\mathcal{F}\subset\Omega$ on which $u$ is not differentiable. 
In this paper, considering the case when $\Omega^* = \Omega^*_1 \cup \Omega^*_2$, with $\Omega^*_1$ and $\Omega^*_2$ being bounded convex domains separated by a hyperplane $\mathcal{H}$, we will establish the regularity of the singular set $\mathcal{F}$.

From previous works \cite{CL1, KM1}, one knows $\mathcal{F} \subset \Omega$ is a Lipschitz hypersurface, partitioning $\Omega$ into two subdomains $\Omega_i$, $i = 1, 2$, such that 
	\begin{equation}\label{condi3} 
		(Du)_{\sharp}(f\chi_{\Omega_i}) = g\chi_{\Omega^*_i}\qquad \text{ for }\quad i = 1, 2.
	\end{equation}
Moreover, $\mathcal{F}$ is characterised as $\partial\Omega_i \cap \Omega$ for $i = 1, 2$, which can be expressed as the graph of a certain Lipschitz function.


Define the functions $u_i : \mathbb{R}^n\rightarrow \mathbb{R}$, for $i = 1, 2$, as follows:
	\begin{equation}\label{uidef}
    		u_i(x) := \sup\{L(x) : L \text{ is affine, } L \leq u \text{ in } \Omega_i, \text{ and } DL(x) \in \Omega_i^*\} \quad \forall x\in \mathbb{R}^n.
	\end{equation}
From \cite{CL1, KM1}, one has $\mathcal{F} = \{u_1 = u_2\} \cap \Omega$. 
Hence, if $u_i$ are differentiable at $x \in \mathcal{F}$, the unit normal vector to the hypersurface $\mathcal{F}$ at $x$ satisfies
\begin{equation}\label{gn1} 
    \nu_\mathcal{F}(x) = \frac{Du_1(x) - Du_2(x)}{|Du_1(x) - Du_2(x)|}.
\end{equation}

\section{Preliminaries}\label{S2} 

Given a convex function $w: \mathbb{R}^n \rightarrow \mathbb{R}$, its subdifferential at $x \in \mathbb{R}^n$ is defined as:
\begin{equation}
\partial^{-}w(x) := \{y \in \mathbb{R}^n : w(z) \geq w(x) + y \cdot (z - x) \ \text{ for all} \ z \in \mathbb{R}^n\}.
\end{equation}
For any subset $A\subset\mathbb{R}^n$, $\partial^{-}w(A):=\cup_{x\in A}\partial^{-}w(x)$.

Let $U \subset \mathbb{R}^n$ be an open set and $C$ be a positive constant. 
We say that a convex function $w$ satisfies the Monge-Ampère inequality:
\begin{equation}
\frac{1}{C} \chi_{_U} \leq \det\, D^2w < C \chi_{_U}
\end{equation}
in the Alexandrov sense if and only if for any Borel set $A$,
\begin{equation}\label{aama}
\frac{1}{C} |A \cap U| \leq |\partial^-w(A)| \leq C |A \cap U|,
\end{equation}
where $|\cdot|$ denotes the $n$-dimensional volume. 

Now,
let us consider the functions and sets $u, v, u_i$ and $\Omega_i, \Omega^*_i$ for $i=1, 2$ as delineated in the introduction. 
Specifically, we have the conditions \eqref{condi1}--\eqref{condi3} and $u_i$ satisfies \eqref{uidef}. 
It is important to note that $\Omega$ can be expressed as the union $\Omega = \Omega_1\cup \Omega_2\cup\mathcal{F}$, and within each $\Omega_i$, $i=1,2$, we have $u_i = u$. 
Additionally, we have the following dual relationships:
\begin{align}
	 v(y) &= \sup_{x\in \Omega}\{y\cdot x - u(x)\} \quad \text{ for all } y\in \mathbb{R}^n; \label{dual111} \\
	 u_i(x) &= \sup_{y\in \Omega_i^*}\{y\cdot x - v(y)\} \quad \text{ for all } x\in \mathbb{R}^n,\ \ i=1,2. \label{dual222}
\end{align} 

Given that $\Omega^*_i$ are convex for $i=1, 2$, Caffarelli's regularity theory \cite{C92} ensures that $u_i\in C^{1,\beta}_{\text{loc}}(\Omega_i)$ for some $\beta\in(0,1)$, and they satisfy the Monge-Ampère inequality in the sense of Alexandrov:
\begin{equation}\label{Asolv} 
    C^{-1}\chi_{\Omega_i} \leq \det\, D^2 u_i \leq C\chi_{\Omega_i} \quad \text{ for } i=1, 2.
\end{equation}
Lastly, we note that $\mathcal{F}=\{x\in\Omega: u_1(x)=u_2(x)\}$.


To proceed further, we introduce some useful definitions and notations.
Let $w:\mathbb{R}^n\rightarrow \mathbb{R}$ be a convex function whose graph does not contain any infinite straight line.

\begin{definition}\label{defS}
Given a point $y_0\in \mathbb{R}^n$ and a small positive constant $h$, we define the centred sections of the convex function $w$ at $y_0$ with height $h$ as:
\begin{equation}\label{sect}
S^c_{h}[w](y_0) := \left\{y\in\mathbb{R}^n : w(y) < w(y_0) + (y-y_0)\cdot \bar{p} + h\right\},
\end{equation}
where $\bar{p}\in \mathbb{R}^n$ is chosen such that the centre of mass of $S^c_{h}[w](y_0)$ coincides with $y_0$.
Additionally, we define the sub-level set of $w$ at $y_0$ with height $h$ as:
\begin{equation}\label{sub}
S_h[w](y_0) := \left\{y\in \mathbb{R}^n : w(y) < \ell_{y_0}(y) + h\right\},
\end{equation}
where $\ell_{y_0}$ denotes an affine supporting function of $w$ at $y_0$.
\end{definition}

\begin{remark}
\emph{For the existence of $\bar{p}$ in \eqref{sect}, we refer the reader to \cite{C92b}.
If $w(0)=0$ and $w\geq 0$, by \cite[Remark 2.5]{CLW3}, we have that 
\begin{equation}\label{secrela}
w \leq Ch \quad \text{in}\ S^c_h[w](0),
\end{equation}
where $C$ is a constant depending only on $n$. }
\end{remark}


In subsequent discussions, a constant will be referred to as universal if it depends solely on the parameters $n, \Omega, \Omega^*, \lambda$, and $\text{dist}(x_0, \partial\Omega)$, where $x_0$ is a generic point of $\mathcal{F}$. Additionally, we will employ the notation $C_n$ to represent a positive constant dependent exclusively on the dimension $n$, with the acknowledgment that its value may change across different contexts. The vectors \( e_i \), where \( i = 1, \ldots, n \), will be used to denote the standard basis vectors in the coordinate system.

\begin{theorem}\label{t111}
Let $\Omega\subset\mathbb{R}^n$ be a bounded domain, and let $\Omega^* = \Omega^*_1 \cup \Omega^*_2$ be a union of bounded convex domains $\Omega^*_1$ and $\Omega^*_2$ in $\mathbb{R}^n$, separated by a hyperplane $\mathcal{H}$. If $\frac{1}{\lambda} < f, g < \lambda$ within $\Omega$ and $\Omega^*$, respectively, for some positive constant $\lambda$, then the hypersurface $\mathcal{F}$ (as long as it is non-empty) is $C^{1,\beta}$ regular inside $\Omega$ for some $\beta \in (0, 1)$. Additionally, $u_i \in C_{\text{loc}}^{1,\beta}(\Omega \cap \overline{\Omega_i})$ for $i = 1, 2$.
\end{theorem}

\section{$C^{1,\beta}$ regularity of $\mathcal{F}.$}\label{S3}

In this section, we continue to utilise the notations established in Section \ref{S2} for $\Omega, \Omega^*, \Omega^*_i$ $(i=1,2)$, $u, v, u_i, \mathcal{F}$. Consider a specific point $x_0\in\mathcal{F} \subset \Omega$. Select sequences $\{x_k\} \subset \Omega_1$ and $\{\hat{x}_k\} \subset \Omega_2$, both converging to $x_0$ as $k \to \infty$. By applying a translation of coordinates, we may assume $x_0 = 0$.
Without loss of generality, and potentially passing to subsequences, we assume $\lim_{k\to\infty}Du(x_k)=y_0 \in \overline{\Omega^*_1}$, and $\lim_{k\to\infty}Du(\hat{x}_k)=\hat{y}_0 \in \overline{\Omega^*_2}$, respectively. This ensures that $y_0, \hat{y}_0 \in \partial^{-}u(0).$ By \eqref{dual111}, we deduce that $0\in\partial^{-}v(y_0) \cap \partial^{-}v(\hat{y}_0)$.

By subtracting an appropriate constant, we may assume $v(y_0)=v(\hat{y}_0)=0$ and $v\geq0$. Define $\Sigma = \{y\in\mathbb{R}^n : v(y) = 0\}$ and denote the set of extreme points of \( \Sigma \) as \( \text{ext}(\Sigma) \).  

Recall Equation \eqref{condi2}, especially $Dv(y) \in \overline{\Omega}$ for almost every $y\in \mathbb{R}^n.$ 
Consider a section of \(v\), namely \(Z:=\{y \in \mathbb{R}^n : v(y) < \ell(y)\}\), with \(\ell\) being an affine function. Provided that \(Z\) is bounded and contains a point \(y\) such that \(0 \in \partial^{-}v(y)\), since \(0\in \mathcal{F}\) is inside the interior of \(\Omega,\)
 we can assert that \(\partial^{-}v(y) \cap \Omega \neq\emptyset\). 
 Denoting the centre of mass of \(Z\) by \(z\), and employing John's lemma, we identify an affine transformation \(L\) such that $L(z)=0$ and \(B_1(0) \subset LZ \subset B_{C_n}(0)\), where \(C_n\) solely depends on the dimension \(n\). 

From \cite[Inequality (4)]{FK1}, we deduce that:
\begin{equation}\label{la1}
|\inf_{Z} (v-\ell)|^n \geq C_1 \left| \left( \frac{1}{2}Z \right)_{z} \cap \Omega^* \right| |Z|,
\end{equation}
where $\left(\frac{1}{2}Z\right)_{z}$ represents the dilation of $Z$ with respect to $z$, and $C_1$ is a constant depending on $n$ and $\lambda$.
From \cite[Proposition 1]{FK1}, we have the following bound:
\begin{equation}\label{ua1}
|v(y) - \ell(y)|^n \leq C_2 |Z \cap \Omega^*||Z| \, \text{dist}\big(L(y), \partial(L(Z))\big) \quad\mbox{for } y\in Z,
\end{equation}
where $C_2$ is a constant depending only on $n$, $\lambda$, and the diameter of $\Omega$.






\begin{lemma}\label{sc111}
The set $\Sigma$ is both bounded and closed, and $\text{\normalfont ext}(\Sigma)\subset \partial\Omega^*$.
\end{lemma}

\begin{proof}
First, we prove that $\Sigma$ is bounded and closed. Assume, to the contrary, that this is not the case. 
By convexity, $\Sigma$ would then contain a half-line of the form $\{q + te : t \geq 0\}$ for some $q \in \Sigma$ and $e \in \mathbb{S}^{n-1}$. For any $y \in \mathbb{R}^n$ and $x \in \partial^{-}v(y)$, by the convexity of $v$ one has 
	$$ (x - 0) \cdot (y - q - te) \geq 0. $$ 
As $t$ tends to infinity, this implies that $x \cdot e \leq 0$. 
Given the arbitrariness of $y$, it follows that 
	$$\partial^{-} v(\mathbb{R}^n) \subseteq \{x \cdot e \leq 0\}.$$ 
Since we have $(Dv)_{\sharp}g = f$, it can be deduced that $\Omega \subseteq \partial^{-} v(\mathbb{R}^n) \subseteq \{x \cdot e \leq 0\},$ which contradicts the fact that $0 \in \Omega$. Thus, $\Sigma$ must be bounded and closed.

Next, we prove that $\text{ext}(\Sigma) \subseteq \partial\Omega^*$. Suppose to the contrary that there exists a point $p \in \text{ext}(\Sigma)$ such that $p$ is either in $\Omega^*$ or in $\mathbb{R}^n \setminus \overline{\Omega^*}$. If $p \in \text{ext}(\Sigma) \cap \Omega^*$, then we have $0 \in \partial^{-}v(p) \cap \Omega$. Referring to Proposition 2 and the subsequent discussion in \cite{FK1}, $v$ is inferred to be $C^1$ and strictly convex in a neighborhood of $p$. Consequently, by duality, the function $u$ must be differentiable at $0$, leading to a contradiction. In the case where $p \in \text{\normalfont ext}(\Sigma) \cap (\mathbb{R}^n \setminus \overline{\Omega^*})$, the application of inequalities \eqref{la1} and \eqref{ua1}, along with the arguments presented in Section 3.1 of \cite{FK1}, also yields a contradiction. Hence, we conclude that $\text{ext}(\Sigma) \subset \partial\Omega^*$.
\end{proof}




Since exposed points are dense in the set of extreme points, we can choose points 
	\begin{equation}\label{exppt}
		p\in\text{ext}(\Sigma) \cap \partial\Omega^*_1 \quad \text{ and }\ \  \hat{p} \in  \text{ext}(\Sigma) \cap \partial\Omega^*_2
	\end{equation} 
 such that both are exposed points of \( \Sigma \). In this context, we establish the following localisation properties.

\begin{lemma}\label{locc1}
Given any \( r_0 > 0 \), there exists an \( h_0 > 0 \) such that for all \( h \leq h_0 \), the inclusion \( S^c_h[v](p) \subset B_{r_0}(p) \) holds,  where $p$ is as in \eqref{exppt}. 
\end{lemma}

\begin{proof}
Since \( p \) is an exposed point of \( \Sigma \), there exists a unit vector $e\in\mathbb{S}^{n-1}$ such that 
	\begin{align*} 
		& \Sigma \subset \{ y : (y - p) \cdot e \leq 0 \}, \\
		& \Sigma \cap \{ y : (y - p) \cdot e = 0 \} = \{ p \}.
	\end{align*} 
Since $S^c_h[v](p)$ is centred at $p$, if $S^c_h[v](p) \cap \{ y : (y - p) \cdot e \geq 0 \} \subset B_r(p) $, then it follows that $S^c_h[v](p) \subset B_{C_n r}(p)$.

Observing that \( v > 0 \) in \( \mathbb{R}^n \setminus \Sigma \) and given any small constant \( r_0>0 \), the convexity and continuity of \( v \) ensure the existence of a constant \( \delta_0 \) such that
\begin{equation}\label{pfloc1}
v > \delta_0 \quad \text{in} \ \ \{ y : (y - p) \cdot e \geq 0 \} \setminus B_{\frac{r_0}{C_n}}(p).
\end{equation}

By \eqref{secrela} we have $v \leq C_n h$ in $S^c_h[v](p)$. Setting \( h_0 = \frac{\delta_0}{C_n} \), for all \( h \leq h_0 \), we then have \( v \leq \delta_0 \) in \( S^c_h[v](p) \). From \eqref{pfloc1}, it follows that
\[ S^c_h[v](p) \cap \{ y : (y - p) \cdot e \geq 0 \} \subset B_{\frac{r_0}{C_n}}(p) \quad \forall\ h \leq h_0. \]
Therefore, we conclude that \( S^c_h[v](p) \subset B_{r_0}(p) \) whenever \( h \leq h_0 \).
\end{proof}


The next lemma establishes the localisation of centred sections near the point $p$. 

\begin{lemma}\label{locc2}
For any given \( \frac{1}{2} > r_0 > 0 \), let \( h_0 \) be the constant specified in Lemma \ref{locc1}. There exists a positive constant \( \delta_0 > 0 \) such that for all points \( y \in \overline{\Omega_1^*} \cap B_{\delta_0}(p) \) and for all \( h \leq h_0 \), the inclusion \( S_h^c[v](y) \subset B_{2r_0}(p) \) holds,  where $p$ is as in \eqref{exppt}. 
\end{lemma}

\begin{proof}
Assume, for the sake of contradiction, that the statement is false. Then, there exist sequences \( \{y_k\} \subset \overline{\Omega^{*}_1} \) converging to \( p \) and \( \{h_k\} \) with \( h_k \leq h_0 \), such that 
\begin{equation}\label{cseg1}
S_{h_k}^c[v](y_k) \cap \left(\mathbb{R}^n \setminus B_{2r_0}(p)\right) \neq \emptyset.
\end{equation}
Given that \( S_{h_k}^c[v](y_k) \) is balanced with respect to \( y_k \), there exists a line segment \( I_k \subset S_{h_k}^c[v](y_k) \) centred at \( y_k \) satisfying \( |I_k| \geq C_n r_0 \).
Without loss of generality, by taking a subsequence if necessary, we can assume that \( h_k \to \bar{h} \leq h_0 \) and \( I_k \to I_\infty \) as \( k \to \infty \), where \( I_\infty \) is a segment centred at \( p \) and satisfying
\begin{equation}\label{cseg2}
|I_\infty| \geq C_n r_0.
\end{equation}

First, we claim that $S_{h_k}^c[v](y_k)$ is uniformly bounded. 
From Lemma \ref{sc111}, $\Sigma$ is compact. 
Owing to the non-negativity and convexity of $v$, we deduce the existence of a large constant \( R > 1 \) satisfying
\[ v(y) > h_0 + \sup_{B_1(p)} v \quad \forall\, y \in \mathbb{R}^n \setminus B_R(p). \]

Suppose \( S_{h_k}^c[v](y_k) = \{y \in \mathbb{R}^n : v(y) < \ell_k\} \), for some affine function \( \ell_k \)  that fulfils \( \ell_k(y_k) = v(y_k) + h_k \). For any given unit vector \( e \in \mathbb{S}^n \), we may assume $\ell_k$ is decreasing in the $e$-direction, otherwise consider $-e$ instead.

Since \( y_k \) converges to \( p\), we may assume that \(y_k\in B_1(p) \) for all $k.$
It can be verified that \( y_k + 3Re \in \mathbb{R}^n \setminus B_R(p) \), leading to the inequality
\[ v(y_k + 3Re) > h_0 + \sup_{B_1(p)} v \geq \ell_k(y_k)\geq \ell_k(y_k+3Re). \]
This implies \( y_k + 3Re \notin S_{h_k}^c[v](y_k) \). 
Since \( S_{h_k}^c[v](y_k) \) is centred at \( y_k \), we deduce \( y_k - 3C_nRe \notin S_{h_k}^c[v](y_k) \). Given that \( e \) is arbitrary, we conclude the claim that 
\begin{equation} 
	S_{h_k}^c[v](y_k) \subset B_{3C_nR}(y_k) \subset B_{3C_nR+1}(p).
\end{equation}

Now, let's consider a subsequence (without changing notation) such that \( S_{h_k}^c[v](y_k) \) converges to a limit convex set \( S_\infty \) in the Hausdorff distance, and $h_k\to\bar h\leq h_0$ as $k\to\infty$. 
Write \( \ell_k = x_k \cdot (y - y_k) + h_k \) with \( x_k \in \mathbb{R}^n \), the global Lipschitz continuity of \( v \) ensures that \( \sup_k |x_k| \leq \|v\|_{\text{Lip}} < \infty \). Consequently, after possibly taking another subsequence, we can assert that \( x_k \to \bar{x} \) as \( k \to \infty \).


In the case when \( \bar{h} > 0 \), it follows that
\[
S_\infty = \left\{ y \in \mathbb{R}^n : v(y) < \bar{x} \cdot (y - p) + \bar{h} \right\} = S^c_{\bar{h}}[v](p).
\]
By \eqref{cseg1}, passing to limit we have
\[
S^c_{\bar{h}}[v](p) \cap \left( \mathbb{R}^n \setminus B_{2r_0}(p) \right) \neq \emptyset,
\]
leading to a contradiction with Lemma \ref{locc1}.

In the case when \( \bar{h} = 0 \), we observe that \( 0 \leq v \leq \bar{x} \cdot (y - p) \) for all \( y \) in \( I_\infty \). Given that \( I_\infty \) is centred at \( p \), it follows that \( v \) must be identically zero on \( I_\infty \), implying that \( I_\infty \subseteq \Sigma \). This contradicts to the fact that $p$ is an exposed point of $\Sigma$.
\end{proof}







Take $r_0$ small so that $B_{2r_0}(p)\cap \Omega_2^*=\emptyset$. Let $h_0, \delta_0$ be as in Lemma \ref{locc1} and \ref{locc2}. Define $\rho_0:=\frac{1}{2}\text{dist}(0,\partial \Omega)$. 
We will show the differentiability of $v$ near $p$ in the following lemma.
\begin{lemma}\label{c1lemma}
There exists a constant $\bar{r}_1<\frac{1}{3}\delta_0$, such that $v$ is differentiable at all $y\in \overline{\Omega_1^*}\cap B_{\bar{r}_1}(p)$, where $p$ is as in \eqref{exppt}. 
\end{lemma}


\begin{proof}

First, we establish the differentiability of \( v \) at the point \( p \).
Before giving the detailed proof, let us outline the main idea. Suppose \( v \) is not differentiable at \( p \). By subtracting an affine function, we may assume that $v$ satisfies \eqref{difs1} for some unit vector $e$ pointing outside $\Omega^*_1$. Consider the section \( S^c_h[v](y_{t_0}) = \{ v < \ell \} \), centred at \( y_{t_0} := p - t_0 e \in\Omega^*_1 \), where $t_0>0$ is chosen small so that $y_{t_0}$ lies sufficiently close to $p$. We have two observations: $(i)$  $p$ is nearly the minimum point of $v-\ell,$ and $(ii)$ After applying a normalisation via the affine transformation \( L \) in \eqref{normdd}, \( L(p) \) lies very close to the boundary of \( L\left(S^c_h[v](y_{t_0})\right) \) as $h\to0$. Using the estimates in \eqref{la1} and \eqref{ua1}, we arrive at a contradiction. This argument resembles Caffarelli's localisation technique and the discussion around equations (1.4) and (1.5) in \cite{GK}, but involves new delicate estimates  as $p\in\partial\Omega^*_1$ lies on the boundary. Below, we provide the detailed proof.  


Consider \( n \) linearly independent unit vectors \( \hat{e}_i \) (may not be orthogonal), where \( i = 1, \ldots, n \). 
By convexity of $v$, it suffices to show that \( v(p + t\hat{e}_i) \) is differentiable at \( t = 0 \) as a convex function of a single variable \( t \), for each \( i = 1, \ldots, n \).

Suppose to the contrary that \( v \) is not differentiable at \( p \). Then there exists a unit vector \( e \), such that \( \{p + te : t < 0\} \cap \Omega_1^* \neq \emptyset\), and \( v(p + te) \) is not differentiable at \( t = 0 \). By subtracting an affine function \( \ell_1 \), we can assume the following behaviour (see Figure 3.1):
\begin{equation}\label{difs1}
\begin{aligned}
v &\geq 0&& \text{on}\ \mathbb{R}^n,\\
v(p + te) &= o(t) && \text{for } t < 0, \\
v(p + te) &= at + o(t) && \text{for } t > 0,
\end{aligned}
\end{equation}
where \( a>0 \) is a constant. 
Note that the target of \( Dv \) is changed to \( \tilde{\Omega} = \Omega - \{D\ell_1\} \), a translation of \( \Omega \).
Given that \( \partial^{-}v(p) \cap \tilde{\Omega} \) is non-empty, the estimates \eqref{la1}--\eqref{ua1} still apply.  

Define \( y_{t_0} := p - t_0e \) for $t_0 \in (0, \frac{1}{2}\delta_0)$ sufficiently small. 
Let's consider the centred section 
	$$S^c_{h}[v](y_{t_0}) = \{v < \ell\},$$ 
where $\ell$ is an affine function satisfying $\ell(y_{t_0}) = v(y_{t_0})+h$. 
Consider the intersection
of the boundary $\partial S^c_{h}[v](y_{t_0})$ with the line $\{p+te : t\in\mathbb{R}\}$, and let
	\[ \partial S^c_{h}[v](y_{t_0}) \cap \big\{p+te : t\in\mathbb{R}\big\} =  \big\{ p - t_1(h)e, \ \ p+t_2(h)e \big\}, \]
where $t_1(h) > t_0$ and $t_2(h) > -t_0$. 
By \cite[Lemma A.8]{CM}, we have that \(S^c_{h}[v](y_{t_0})\) varies continuously with respect to \(h \in (0, h_0]\). Hence, \(t_1(h)\) and \(t_2(h)\) depend continuously on \(h\).
Since \(S^c_{h}[v](y_{t_0})\) is centred at $y_{t_0}$, we have the following bounds:
\begin{equation}\label{t0t1}
t_0 < t_1(h) < C_n(t_0 + t_2(h)),
\end{equation}
where \(C_n\) is a constant depending only on \(n\).


\begin{figure}\label{f3.1}
 \centering
 \includegraphics[scale=0.2]{2extrem.png}
 \caption{A centred section \(S^c_{h}[v](y_{t_0})\), near $p$. }
\end{figure}

We first \emph{claim} that \(t_2(h_0) \geq t_0\), provided $t_0$ is chosen sufficiently small initially.
Suppose, to the contrary, that \(t_2(h_0) < t_0\).
Then, by \eqref{difs1}, we have 
\begin{equation}\label{t21}
\ell(p+t_2(h_0)e) = v(p+t_2(h_0)e) \leq 2at_0,
\end{equation}
provided \(t_0\) is small enough.
By \eqref{difs1} and \eqref{t0t1}, we obtain 
\begin{equation}\label{t22}
\ell(p-t_1(h_0)e) = v(p-t_1(h_0)e) = o(t_0).
\end{equation}
By \eqref{t21}, \eqref{t22}, and the balance of \(S^c_{h_0}[v](y_{t_0})\) with respect to $y_{t_0}$,
we have 
	\[ \ell(y_{t_0}) \leq C_n(2at_0 + o(t_0)), \] contradicting the fact that \(\ell(y_{t_0}) \geq \ell(y_{t_0}) - v(y_{t_0}) = h_0\), provided \(t_0 \ll h_0\) was chosen initially.
Hence, the claim is proved.


We then \emph{claim} that \(\liminf_{h\rightarrow 0}t_2(h)\leq 0\), provided $t_0$ is chosen sufficiently small initially. 
Suppose, to the contrary, that there exists a constant \(a_0 \in (0, t_0)\) and a sequence \(\{h_j\}\) converging to \(0\), such that \(t_2(h_j) \geq a_0\) for all \(j\). Hence, from \eqref{difs1},
\begin{equation}\label{t2hj}
v(p+t_2(h_j)e)\geq \frac{1}{2}aa_0,
\end{equation}
provided \(t_0\) is small enough.
Moreover, since $S^c_{h_j}[v](y_{t_0})$ is centred at $y_{t_0}$ we have the segment 
\begin{equation}
I=\left\{p+\theta e: -t_0-\frac{1}{C_n}t_0<\theta<a_0\right\} \subset S^c_{h_j}[v](y_{t_0})\quad\forall\,j=1,2,\cdots.
\end{equation}
Suppose \(S^c_{h_j}[v](y_{t_0})=\{y\in\mathbb{R}^n: v(y)<v(y_{t_0})+x_j\cdot (y-y_{t_0})+h_j\}\) for some vector \(x_j\in \mathbb{R}^n\).
Since \(v\) is globally Lipschitz, \( \sup_j |x_j| \leq \|v\|_{\text{Lip}} < \infty \). Consequently, by taking a subsequence if necessary, we may assume \(x_j \to \bar{x}\) as \(j \to \infty\).
In the limit, we thus obtain 
	\[ 0 \leq v(y) \leq \bar{\ell}(y) \quad \forall\, y \in I, \] 
where \( \bar{\ell}(y) = v(y_{t_0}) + \bar{x} \cdot (y - y_{t_0}) \). 
Since \(y_{t_0} \in I\), \(v(y_{t_0}) = \bar{\ell}(y_{t_0})\) and \(v \leq \bar{\ell}\) on \(I\), the convexity of \(v\) implies that \(v \equiv \bar{\ell}\) on \(I\), namely \(v(p + te)\) is affine with respect to \(t\) for \(t \in (-t_0 - \frac{1}{C_n}t_0, a_0)\), which contradicts \eqref{difs1}. 
Hence, the claim $\liminf_{h \to 0} t_2(h) \leq 0$ is proved.



Now, choose a $t_0>0$ small enough such that the above two claims hold. By these claims and the continuity of \(t_2(h)\) with respect to \(h\), there exists an \(h'\in(0, h_0)\) such that 
\begin{equation}
t_2(h') = \frac{1}{M}t_0,
\end{equation}
where \(M > 1\) is a large constant to be determined.
By \eqref{t0t1}, we then have 
\begin{equation}
t_0 < t_1(h') < C_n(t_0 + t_2(h')) \leq 2C_nt_0.
\end{equation}
Assume that \(S^c_{h'}[v](y_{t_0}) = \{v < \ell\}\) for some affine function \(\ell\) satisfying \(\ell(y_{t_0}) = v(y_{t_0}) + h'\).


Evaluating the affine function \(\ell\) at specific points yields:
\begin{align}
\ell(p - t_1e) = v(p - t_1(h')e) &= o(t_1) = o(t_0), \\
\ell\left(p + t_2(h')e\right) = v\left(p + \frac{1}{M}t_0e\right) &= \frac{a}{M}t_0 + o(t_0).
\end{align}
The above equations imply that \(\ell\) is increasing along the \(e\) direction (provided $t_0$ is chosen sufficiently small at the beginning), leading to:
\begin{equation}\label{makec1}
\ell(p) - v(p) = \ell(p) \geq \ell(y_{t_0}) \geq \ell(y_{t_0}) - v(y_{t_0}) \approx \left| \inf_{S^c_{h'}[v](y_{t_0})} (v - \ell) \right|.
\end{equation}

Applying John's lemma, we identify an affine transformation \(L\) such that \(L(y_{t_0}) = 0\), and 
\begin{equation}\label{normdd}
B_1(0) \subset L\left(S^c_{h'}[v](y_{t_0})\right) \subset B_{C_n}(0).
\end{equation}
Since the affine transformation preserves the ratio of distances, we have
	\[ \frac{\text{dist}\left(L(p), \partial L\left(S^c_{h'}[v](y_{t_0})\right)\right) }{\text{diam}(L\left(S^c_{h'}[v](y_{t_0})\right))} = \frac{\text{dist}(p,\partial S^c_{h'}[v](y_{t_0})}{\text{diam}(S^c_{h'}[v](y_{t_0})} \leq \frac{t_2(h')}{t_1(h')} \leq \frac{1}{M}.\]
Hence, by \eqref{normdd}
\begin{equation}\label{makec2}
\text{dist}\left(L(p), \partial L\left(S^c_{h'}[v](y_{t_0})\right)\right)   \leq \frac{C_n}{M}.
\end{equation}

Since \(y_{t_0} \in \Omega_1^* \cap B_{\delta_0}(p)\) and \(h' \leq h_0\), Lemma \ref{locc2} implies \(S^c_{h'}[v](y_{t_0}) \subset B_{2r_0}(p)\), ensuring that \(S^c_{h'}[v](y_{t_0}) \cap \Omega_2^* = \emptyset\). 
The convexity of $\Omega^*_1$ leads to the volume comparison:
\begin{equation}\label{makec3}
\left| \left(\frac{1}{2}S^c_{h'}[v](y_{t_0})\right)_{y_{t_0}} \cap \Omega^* \right| \geq C_n \left| S^c_{h'}[v](y_{t_0}) \cap \Omega^* \right|.
\end{equation}
Then by the estimate \eqref{la1}, we can obtain
\begin{equation}\label{makec4}
 \left|\inf_{S^c_{h'}[v](y_{t_0})} (v-\ell)\right|^n \geq C_1 C_n |S^c_{h'}[v](y_{t_0}) \cap \Omega^*| |S^c_{h'}[v](y_{t_0})|,
\end{equation}
where \(C_1\) and \(C_n\) are constants independent of \( M \).

On the other hand, by estimates \eqref{ua1}, \eqref{makec1} and \eqref{makec2}, we can deduce
\begin{equation}
 \left|\inf_{S^c_{h'}[v](y_{t_0})} (v-\ell)\right|^n \lesssim |\ell(p)-v(p)|^n \leq \frac{C_2 C_n}{M} |S^c_{h'}[v](y_{t_0}) \cap \Omega^*| |S^c_{h'}[v](y_{t_0})|,
\end{equation}
where \(C_2\) is another constant  independent of \( M \).
However, this contradicts the estimate \eqref{makec4}, when the constant \(M\) is sufficiently large. 

Consequently, we conclude that the function \(v\) is differentiable at the point \(p\). By selecting a sufficiently small \(\bar r_1 < \delta_0\), we further obtain:
\begin{equation}\label{newextend1}
\partial^{-}v(B_{\bar r_1}(p)) \subset B_{\rho_0}(0) \subset \Omega.
\end{equation}
Last, for any point \(y \in \overline{\Omega_1^*} \cap B_{\bar r_1}(p)\), we can employ the previous argumentation for establishing the differentiability of \(v\) at \(p\) to similarly show that \(v\) is differentiable at \(y\).
\end{proof}

\begin{remark}
\emph{ Let \( r_0, \delta_0, \rho_0, \) and \( \bar r_1 \) be the constants given in preceding lemmas. 
According to \eqref{newextend1}, for any point \( y \in \Omega_1^* \cap B_{\bar r_1}(p) \), the set \( \partial^{-}v(y) \cap \Omega \) is non-empty. By invoking Proposition 2 from \cite{FK1} and the subsequent discussions, it follows that the function \( v \) is continuously differentiable and exhibits strict convexity within \( \Omega_1^* \cap B_{\bar r_1}(p) \). }
\end{remark}


\begin{lemma}[Localisation]
The function $v$ satisfies the Monge-Amp\`ere inequality:
\begin{equation}\label{lockey1}
\frac{1}{\lambda^2}\chi_{_{\Omega^*_1}} \leq \det\, D^2v \leq \lambda^2\chi_{_{\Omega^*_1}} 
\quad \text{in} \ B_{\bar r_1}(p)
\end{equation}
in the Alexandrov sense, as in \eqref{aama}. 
\end{lemma}

\begin{proof}
It suffices to show $\det D^2v$ has no mass in $B_{\bar r_1}(p) \setminus \Omega_1^*$, namely $|\partial^{-}v(B_{\bar r_1}(p) \setminus \Omega_1^*)| =0$. 

In fact, let $ y\in B_{\bar r_1}(p) \setminus \Omega_1^*$, and $x\in \partial^{-}v(y)$. Owing to \eqref{newextend1}, it is guaranteed that \( x \) belongs to \( \Omega = \Omega_1 \cup \Omega_2 \cup \mathcal{F} \).
Recall that $u_i$ is $C^1$ and strictly convex inside $\Omega_i$ for $i=1, 2.$
 If \( x \) were to lie in either of \( \Omega_i \), it would imply that \( y \in\Omega_i^* \), being the gradient \( Du_i(x) \). This, however, contradicts our initial assumption that \( y \in B_{\bar r_1}(p) \setminus \Omega_1^* \). Consequently, it must be that \( x \) resides in the set \( \mathcal{F} \).

This reasoning concludes that the subdifferential \( \partial^{-}v \) maps \( B_{\bar r_1}(p) \setminus \Omega_1^* \) entirely into \( \mathcal{F} \). 
 By \cite[Theorem 5.1]{CL1}, we have that $\mathcal{F}$ is a Lipschitz hypersurface, and hence its Lebesgue measure is zero.
Therefore, the Lebesgue measure of the image set satisfies $|\partial^{-}v(B_{\bar r_1}(p) \setminus \Omega_1^*)| =0$.
\end{proof}



Utilising Lemmas \ref{locc1} and \ref{locc2}, we are able to determine constants \( \bar r_0 < \bar r_1 \) and \( h_0 > 0 \) such that:
\begin{equation}\label{lockey2}
S_h^c[v](y) \subset B_{\bar r_1}(p)  \quad \forall\, y \in B_{\bar r_0}(p) \cap \Omega_1^*, \ \  \forall\, h \leq h_0.
\end{equation}
With these preparations, we are now ready to prove Theorem \ref{t111} as follows.
\begin{proof}[Proof of Theorem \ref{t111}]
By \eqref{lockey1} and \eqref{lockey2}, we can invoke the proof of \cite[Lemma 2.2, Corollary 2.3]{C96} (see also \cite[Theorem 7.13]{CM}) to obtain a quantitative uniform convexity of $v$ within $B_{\bar{r}_0}(p)\cap \overline{\Omega_1^*}$. Specifically, this implies that for any two points $y, \tilde{y} \in B_{\bar{r}_0}(p) \cap \overline{\Omega_1^*}$, one has the inequality 
	$$ |Dv(y) - Dv(\tilde{y})| \geq C|y - \tilde{y}|^a,$$ 
where $a > 2$ and $C$ is a positive constant.

Since $u_1 = v^*$ in $\Omega_1$, where $v^*$ is the Legendre transform of $v$, a basic and well known property of the Legendre transform \cite[Remark 7.10]{CM} allows us to deduce that $u_1 \in C^{1,\beta}(B_r(0) \cap \overline{\Omega_1})$ for some $\beta \in (0, 1)$ and $r > 0$. 
By a similar argument, and a possibly smaller $r$, we can ensure that $u_2 \in C^{1,\beta}(B_r(0) \cap \overline{\Omega_2})$.

Inferred from \eqref{gn1}, we conclude that the singular set $\mathcal{F}$ exhibits $C^{1,\beta}$ regularity in $B_r(0)$. Finally, a standard covering argument ensures the $C^{1,\beta}$ regularity of $\mathcal{F}$ inside $\Omega$.
\end{proof}

\begin{remark}
\emph{
Since \( u_i \in C^{1,\beta}(B_r(0) \cap \overline{\Omega_i}) \) for \( i = 1, 2 \), we have \( p = y_0 = Du_1(0) \) and \( \hat{p} = \hat{y}_0 = Du_2(0)\).
}
\end{remark}

\begin{thebibliography}{12}
 

\bibitem{Br1}
Y. Brenier, 
Polar factorization and monotone rearrangement of vector-valued functions. 
\emph{Comm. Pure Appl. Math.} 44 (1991), no. 4, 375--417.
 
\bibitem{C92}
 L. A. Caffarelli, 
 The regularity of mappings with a convex potential. 
 \emph{J. Amer. Math. Soc.} 5 (1992), 99--104.
 
\bibitem{C92b}
 L. A. Caffarelli, 
 Boundary regularity of maps with convex potentials. 
 \emph{Comm. Pure Appl. Math.} 45 (1992), 1141--1151.

\bibitem{C96}
L. A. Caffarelli, 
Boundary regularity of maps with convex potentials. II.
 \emph{Ann. of Math.} 144 (1996), 453--496.

 \bibitem{CM}  
 L. A. Caffarelli and R. J. McCann, 
 Free boundaries in optimal transport and Monge-Amp\`{e}re obstacle problems.
  \textit{Ann. of Math.} 171 (2010), 673--730.
  
\bibitem{CLL} 
S. Chen, Y. Li and J. Liu,
	Optimal (partial) transport between non-convex polygonal domains.
	Preprint, available at arXiv:2311.15655.

\bibitem{CL1}
S. Chen and J. Liu, 
Regularity of free boundaries in optimal transportation. 
Preprint, available at arXiv:1911.10790. 

\bibitem{CL2}
S. Chen and J. Liu, 
Global regularity in the Monge-Amp\`ere obstacle problem. 
Preprint, available arXiv:2307.00262.

\bibitem{CLW1} 
S. Chen, J. Liu and X.-J. Wang, 
Global regularity for the Monge-Amp\`ere equation with natural boundary condition.    
\textit{Ann. of Math.} (2) 194 (2021), no.3, 745--793.
 

\bibitem{CLW3}
S. Chen, J. Liu and X.-J. Wang,
$C^{2,\alpha}$ regularity of free boundaries in optimal transportation.  
\emph{Comm. Pure Appl. Math.} 77 (2024), 731--794. 

%\bibitem{DeFi}
%G. De Philippis and A. Figalli,
%The Monge-Amp\`ere equation and its link to optimal transportation.
%\emph{Bull. Amer. Math. Soc.} 51 (2014), 527--580.

\bibitem{D91} 
Ph. Delano\"e,
	    Classical solvability in dimension two of the second boundary value problem associated with the Monge-Amp\`ere operator.
	    \textit{Ann. Inst. Henri Poincar\'e, Analyse Non Lin\'eaire} 8 (1991), 443--457.

\bibitem{AF09} 
A. Figalli, 
A note on the regularity of the free boundaries in the optimal partial transport problem.
 \textit{Rend. Circ. Mat. Palermo} 58 (2009), no. 2, 283--286.
 	    
 \bibitem{AFi} 
 A. Figalli, 
 The optimal partial transport problem. 
  \textit{Arch. Ration. Mech. Anal.} 195 (2010), 533--560.
  
 \bibitem{FK1}  
 A. Figalli and Y.-H. Kim, 
 Partial regularity of Brenier solutions of the Monge–Amp\`ere equation. 
 \textit{Discrete Contin. Dyn. Syst.} 28 (2010), 559--565.

 \bibitem{GK}
 N. Guillen and J. Kitagawa,
 On the local geometry of maps with c-convex potentials.
 \textit{Calc. Var. Partial Differential Equations} 52 (2015), 345--387. 

 \bibitem{I} E. Indrei, Free boundary regularity in the optimal partial transport problem. 
 \textit{J. Funct. Anal.} 264 (2013), no. 11, 2497--2528.   
 

\bibitem{KM1} 
J. Kitagawa and R. McCann,
 Free discontinuities in optimal transport.
       \emph{Arch. Rational Mech. Anal.} 232 (2019), 1505--1541. 


\bibitem{SY}
   O. Savin and H. Yu,
    Regularity of optimal transport between planar convex domains. 
    \textit{Duke Math. J.} 169 (2020), 1305--1327.
   
   
\bibitem{U1} J. Urbas,
    On the second boundary value problem of Monge-Amp\`ere type.
      \textit{J. Reine Angew. Math.}  487 (1997), 115--124.


\bibitem{V1} C. Villani,
           \emph{Topics in optimal transportation}.  Grad. Stud. Math. 58,
          \textit{Amer. Math. Soc.}, 2003.

\bibitem{V2} C. Villani,
           \emph{Optimal transport, old and new}.  
           Springer, Berlin, 2006.

\end{thebibliography}
\end{document}
